\section{R ile çözümlü uygulama}
\label{sec:r-b13}

\textbf{Soru.} Bölümdeki 16 gözlemlik çalışma saati--son test örneğini
R'ye aktarın. Pearson korelasyonunu, regresyon doğrusunu, $R^2$'yi ve altı
saat için öngörüyü bulun; artık grafiklerini inceleyin.

\begin{lstlisting}[style=prismR]
veri <- data.frame(
  saat = c(5, 3, 7, 6, 6, 2, 8, 4, 4, 5, 3, 7, 3, 5, 2, 6),
  puan = c(60, 54, 68, 63, 65, 52, 70, 58,
           57, 61, 52, 64, 55, 60, 50, 63)
)
cor.test(veri$saat, veri$puan, method = "pearson")
model <- lm(puan ~ saat, data = veri)
summary(model)
confint(model)
yeni <- data.frame(saat = 6)
predict(model, newdata = yeni, interval = "confidence")
predict(model, newdata = yeni, interval = "prediction")
plot(puan ~ saat, data = veri)
abline(model, col = "blue")
plot(fitted(model), resid(model),
     xlab = "Uydurulan deger", ylab = "Artik")
abline(h = 0, lty = 2)
qqnorm(resid(model))
qqline(resid(model))
\end{lstlisting}

\begin{outputguidebox}
\textbf{Çözüm ve kontrol değerleri.} $r=0{,}9806$ ve
$R^2=0{,}9616$'dır. Regresyon doğrusu
$\hat y=44{,}5980+3{,}1373\,\text{saat}$ olur. Altı saat için nokta
öngörüsü 63,4216 puandır. Eğimin sıfır olduğu hipotezinin çift yönlü
\pval-değeri yaklaşık $2{,}62\times10^{-11}$'dir; bu sayı sıfır değildir.
\end{outputguidebox}

\textbf{Yorum.} \texttt{confidence} aynı saatteki evren ortalaması için,
\texttt{prediction} ise yeni bir bireyin puanı için aralık verir. İkincisi
bireysel değişkenliği de içerdiğinden daha geniştir. Doğrusallık, sabit hata
varyansı, bağımsızlık ve çıkarım için hata dağılımı değerlendirilmelidir.

\textbf{Pekiştirme ve yanıt.} Altı saat gözlenen 2--8 saat aralığındadır;
20 saat için öngörü dışa taşırmadır. Güçlü korelasyon, ek bir saat çalışmanın
her öğrencinin puanını nedensel olarak 3,1373 artırdığını göstermez.